\documentclass[11pt,a4paper]{article}
\usepackage[margin=19mm,headheight=15pt]{geometry}
\usepackage{fontspec}
\setmainfont{Noto Serif}
\usepackage{amsmath,amssymb,mathtools}
\usepackage{xcolor}
\usepackage{fancyhdr}
\usepackage{enumitem}
\usepackage{microtype}
\setlength{\parindent}{0pt}
\setlength{\parskip}{6pt}
\definecolor{navy}{HTML}{183153}
\definecolor{lightgray}{HTML}{F3F5F7}
\pagestyle{fancy}\fancyhf{}
\lhead{\small ME 5309 -- Mid-Semester Examination}
\rhead{\small ME 5309}
\cfoot{\thepage}
\begin{document}
\begin{center}
{\Large\bfseries ME 5309 -- Mid-Semester Examination}\\[3pt]
{\large Answer Booklet}\\[8pt]
\end{center}
\textbf{Instructions:} Attempt any three questions. Each question carries 10 marks. Show working where appropriate.
\section*{Question 1 \hfill [10 marks]}
\subsection*{Q1(a) \hfill [4 marks]}

\vspace{2mm}

\subsection*{Q1(b) \hfill [3 marks]}

\vspace{2mm}

\subsection*{Q1(c) \hfill [3 marks]}

\vspace{2mm}

\section*{Question 2 \hfill [10 marks]}
\subsection*{Q2(a) \hfill [2 marks]}

The state model is ẋ₁ = x₂, ẋ₂ = (τ-0.05x₂-2.4525 sin x₁)/0.02.
Equilibrium requires τ₀ = 2.4525 sin θ\_e, so τ₀ must be between -2.4525 and +2.4525 N m.

\subsection*{Q2(b) \hfill [2 marks]}

\[\begin{aligned}A = [[0,1], [-61.3125,-2.5]], B = [[0],[50]]\end{aligned}\]

\subsection*{Q2(c) \hfill [4 marks]}

The balance is 50(90-θss) = 2.4525 sin θss.
Solving gives a small error of about 0.05°.

\subsection*{Q2(d) \hfill [2 marks]}

The added gravity torque should cancel the load, so the steady-state error becomes zero. 
\section*{Question 3 \hfill [10 marks]}
\subsection*{Q3(a) \hfill [2 marks]}

The estimate is noisy and can fluctuate from sample to sample.

\subsection*{Q3(b) \hfill [2 marks]}

\[\begin{aligned}y[k]=0.8818y[k-1]+0.0591x_k+0.0591x[k-1], DC gain 1\end{aligned}\]

\subsection*{Q3(c) \hfill [4 marks]}

\vspace{2mm}

\subsection*{Q3(d) \hfill [2 marks]}

The observer is simpler because it only smooths the measured velocity.
The filter is safer when the model is poor.

\section*{Question 4 \hfill [10 marks]}
\subsection*{Q4(a) \hfill [2 marks]}

The closed-loop poles are -2±2j.
I use Nff=c/b=6 for unity DC gain.

\subsection*{Q4(b) \hfill [2 marks]}

\[\begin{aligned}\theta_{ss} = 0.75\times8/(6+2)=0.75,\\ \text{the percentage tracking error is 25\%}\end{aligned}\]

\subsection*{Q4(c) \hfill [4 marks]}

\vspace{2mm}

\subsection*{Q4(d) \hfill [2 marks]}

Integral action gives zero steady-state error if stable.
A practical problem is integrator windup.

\end{document}